\documentclass[12pt]{article}

\usepackage[a4paper, total={6in, 8in}]{geometry}
\usepackage{textcomp}

\begin{document}
\setlength{\parindent}{0pt}

\def \t {\textrightarrow}
\def \v {\vspace{1em}}
\def \bi {\begin{itemize}}
\def \ei {\end{itemize}}
\def \s[#1] {\section*{#1}}
\def \ss[#1] {\subsection*{#1}}
\def \sss[#1] {\subsubsection*{#1}}

\s[Recuperare 11 maggio]

\s[UE - Storia]
Vedere "La storia dell'UE in pillole" \\
UE nasce con il manifesto di Ventotene (scritto di Spinelli, Rossi etc.) \t è la base dell europa \\
Nel 49 nasce il consiglio d'europa \t che nasce a tutela dei diritti umani \t deve garantire omogeneita nella difesa dei diritti umani \\
Questo è il primo tentativo europeo di creare un legame tra i paesi che escono dalla guerra \\
Nel 50 Schumann è ministro estero della Francia \t e propone la creazione della CECA \t l'idea dell'unione nasce economica \\
Es. Zwolferein \t piano economico prima della germania unita, proposta un'unione economica doganale \t la ricchezza comune e il passaggio di merci pone infatti i presupposti \\
La CECA doveva garantire legami economici, in modo da rendere inpensabile una guerra tra stati europei \t l economia garantisce la pace perche non si attaccano partner finanziari

\v

Nel 1951 alcuni paesi, che hanno istituito la CECA, pensano a una autorita per garantire questo trattato \t si pensa a un alta autorita con sede a Lussenburgo, che doveva essere potere sovranazionale che faceva rispettare le regole del mercato \\
Questi internventi hanno grande successo \t crescita economica, in clima difficile della ricostruzione post bellica \t cosi viene fondata la CEE nel 1957 (l'atto fondativo formato a roma) \\
Nasce poi anche l Euratom (per l'energia atomica, per discorso energetico e quindi economico-commerciale) \\
Nel 57 i trattati di Roma vengono firmati, e entrano in vigore nel 58 \t vengono istituite quindi le commissioni che devono sovraintendere al funzionamento della cee e dell euratom

\v

Nel 68 nasce l unione doganale \t sempre con orientamento economico, si pensa di abolire i dazi \t unione doganale nasce tra i 6 paesi fondatori \t + tariffe commerciali comuni per il mercato estero \\
Gia dal 68 inizia a prendere piede l'idea che serve una moneta comune \t per permette circolazione delle merci e unione economica è + utile \\
L idea inizia gia qua, euro poi sara molto piu tardi \\
Si pensa a questo pero perche le monete hanno valori finanziari molto diversi \t che fluttuano, e questo non garantisce stabilita al mercato \t con moneta unica invece il mercato sarbbe + stabile + redditizio

\v

Intanto la cee cresce, perche altri paesi chiedono di entrare (come nel 73, che entra anche UK e diventano 9 i paesi) \\
Poi si inizia a pensare anche che questo consiglio europeo non sia + sufficiente da solo \t si pensa quindi a degli organismi di potere che alzino il livello della cooperazione \\
Si pensa che l'unione puo diventare anche politica (e non solo economica) \t questo è il grande obbiettivo dell UE, che non è ancora stata raggiunta completamente ora \\
Si pensa quindi a un parlamento europeo, ovveor un organismo che comprenda i rappresentanti dei paesi \t si pensa anche un passaporto europeo \\
Nel 75 il consiglio europeo che prende queste decisioni è presieduto da Aldo Moro \t anche De Gasperi era stato sostenitore di UE \\
Nel 79 si tengono elezioni quindi a suffragio universale

\v

Nel 1984 parlamento UE approva istituzione dell UE \t sostenuto da Spinelli \\
Differenza tra UE e CEE sta nel passaggio da identita solo economica a una anche politica \t sono due cose diverse, anche nella loro denominazione \t c'è un salto di livello

\v

Nel 85 viene firmato Schengen \t è il primo atto ufficiale formale di quello che portera a nascita UE \t stesso anno stilato anche "L'atto unico europeo" = atto che doveva completare il disegno di questa UE, ancora con grande impronta econmica ma andando verso discorso politico \\
Nasce con Schengen un confine territoriale inoltre \t non a Maastricht \t nasce area Schengen

\v

Nel 87 nasce l'Erasmus \t regola passaggio degli studenti di stato in stato, cosi che gli studenti possono essere studenti della comunita europea \\
È un primo tentativo di allargare la base economica della cee \t non si parla + solo di merci e di soldi, ma anche di possibilita per i giovani \\
Non è tappa fondamentale, pero ha rilevanza ideologica \t si inizia ad allargare il discorso

\v

Nel 92 viene firmato il trattato di Maastricht \t ufficializzato il passaggio da CEE a UE \t e vengono firmati diversi trattati che definiscono questo passaggio (politica interna, giustiza, moneta etc.) \\
L UE non è quindi soltanto una CEE che cambia nome \t è qualcosa di diverso: ha alla base un mercato comune, ma UE poi coinvolge gli stati anche in altri ambiti come la sicurezza, linea della politica estera, etc. \\
C'è quindi un salto qualitativo \\
Entra poi in vigore il mercato libero \t circolano liberamente capitali, persone, servizi, merci etc. \\
Nel 93 entra in vigore mercato libero, mentre gli stati membri aumentao notevolmete

\v

Nel 97 trattato di Amsterdam \t tocca relazione tra il cittadino e UE \t in temi di sicurezza, la sua liberta nell UE, etc. \\
È un trattato + particolare, che si occupa dei diritti e doveri del cittadino europeo

\v

Nel 2002 arriva l'euro, e le banconote diventano uguali in tutti gli stati \t c'è anche trattativa degli stati per il cambio in questa valuta, e alcuni paesi vengono sfavoreggiati \\
Nel 2004 per motivi molto finanziari entrano altri paesi \t i paesi fuori penalizzano quelli fuori, perche moneta unica è piu forte \\
Nel 2004 si ha estensione + grande \\
C'è pero allargamento dell UE anche a fronte di cambiamenti politici nel mondo (come caduta urss etc.) \\
Nel 2007 viene firmato anche il trattato di Lisbona, che ridisegna la struttura istituzionale dell UE \t introduce il presidente del consiglio europeo e altor rappresentante \\
Inoltre rafforza potere del parlamento, e fa in modo che i cittadini degli stati possano essere + partecipi nella vita politica eu \t si tenta di dare maggiore stabilita politica e avvicinare i cittadini \\
Fino al 2007 infatti l UE viene vissuta da alcuni cittadini come qualcosa di lontano \\
Nel 2012 ue ottiene nobel per la pace \t si capisce la ratio dietro questo organo che è nato \t ue vuole tutelare i dritti, la pace, etc

\v

2014 \t elezioni europee, e italia assume presidenza del consiglio della UE per 6 mesi \t in questo momento il problema da affrontare è quella dei migranti (strage lampedusa)

\s[Italia nel secondo dopo guerra]
\bi
  \item fine della guerra, primo governo Parri che dura 6 mesi (ex partigiano) 
  \item nel 45 De Gasperi, fino al 53
  \item nel 46 De Nicola capo provissorio, firma cost.
  \item 47 trattati di pace con conferenza di Londra \t e anno della dottrina Truma e viaggio di De Gasperi in USA \t inizia centrismo
  \item 48 referendum per repubblica, prime elezioni \t d.c. stra vince
  \item nel 49 nasce la NATO
  \item 1958 italia entra nella CEE \t questo da impulso al miracolo economico
  \item crisi progressiva centrismo, tra 62-62 e fino 68 centro sinistra
  \item 68 inizia contestazione
  \item 69 autunno caldo, e piazza fontana
  \item compromesso storico del 73 \t Berlinguer fa proposta a Aldo Moro per trovare punto di unione per risolvere terrorismo
  \item 78 rapimento Moro, con questo rapimento inizia il declino delle brigate rosse
  \item primi anni 80 è nuova era, che è quella del penta partito = 5 partiti di maggioranza, con Spadolini capo del govenro (che è repubblicano, primo govenro non democristiano)
  \item anni 80 c'è massoneria, lotta alla mafia, Falcone e Borsellino
  \item 1992 Tangentopoli, che segna la fine della prima repubblica (guidata dai grandi partiti di massa, che non esisteranno + dopo questo)
\ei
\end{document}
